\chapter{Metodolog\'ia de desarrollo e implementaci\'on}\label{ch:metodologia}

\begin{contenidoNuevo}

Este cap\'itulo documenta las decisiones tecnol\'ogicas, la arquitectura de datos, los diagramas de dise\~no y el avance funcional demostrado en el hito del 50\,\%. El texto nuevo se presenta en azul.

\section{Metodolog\'ia}

El desarrollo sigui\'o un enfoque iterativo:
\begin{enumerate}
	\item ingenier\'ia de datos sobre la BD core de la ART (vistas y panel establecimiento$\times$mes);
	\item entrenamiento y evaluaci\'on de modelos XGBoost (horizontes 1m y 3m);
	\item exportaci\'on de ranking explicable (SHAP) e importaci\'on a BD operativa;
	\item construcci\'on de API REST + frontend React para uso por prevencionistas;
	\item anal\'itica hist\'orica (snapshots, cobertura y precisi\'on observada).
\end{enumerate}

La validaci\'on del modelo usa split temporal: entrenamiento 2022--2023, calibraci\'on 2024 y holdout 2025.

\section{Stack tecnol\'ogico y justificaci\'on}

\begin{table}[htbp]
	\centering
	\caption{Stack tecnol\'ogico del prototipo y motivaci\'on de la elecci\'on.}
	\label{tab:stack}
	\begin{tabular}{@{}p{0.28\textwidth}p{0.66\textwidth}@{}}
		\hline
		\textbf{Componente} & \textbf{Elecci\'on y justificaci\'on} \\
		\hline
		Pipeline ML & Python + XGBoost + SHAP: adecuado a datos tabulares heterog\'eneos y a la necesidad de explicabilidad. \\
		API & ASP.NET Core (C\#): encaje con ecosistemas ART/.NET, Minimal APIs y JWT. \\
		Frontend & React + Vite + TypeScript: UI moderna, tipado y ciclo de desarrollo \'agil. \\
		BD operativa & SQL Server Express (`PfiPrevencion`): modelo relacional de rankings, agenda y usuarios. \\
		Auth & JWT Bearer + hash de contrase\~nas: simple para tesis y upgradeable. \\
		Comunicaci\'on & HTTPS/HTTP + REST JSON: desacople cliente-servidor. \\
		Cloud objetivo & AWS Free Tier (EC2 todo-en-uno + S3 + Budgets): bajo costo para demo posterior. \\
		\hline
	\end{tabular}
\end{table}

\subsection{Protocolos y arquitectura de red}

\begin{itemize}
	\item \textbf{REST/HTTP(S):} intercambio de recursos JSON entre React y la API.
	\item \textbf{JWT Bearer:} autorizaci\'on por token en el encabezado \texttt{Authorization}.
	\item \textbf{SQL/TDS:} persistencia en la BD operativa local (y futura en EC2).
	\item \textbf{Separaci\'on de redes:} la BD core de la ART permanece on-prem; solo el pipeline (VPN/local) la consulta en solo lectura.
\end{itemize}

\section{Flujo de informaci\'on}

\begin{figure}[htbp]
	\centering
	% Flujo de informaci\'on del sistema
\begin{tikzpicture}[
  font=\scriptsize\sffamily,
  node distance=0.55cm and 0.7cm,
  box/.style={draw=blue!70, fill=blue!6, rounded corners=2pt, align=center, minimum width=2.3cm, minimum height=0.85cm},
  arr/.style={-{Latex}, thick, blue!60}
  ]
  \node[box] (core) {BD core\\ART (on-prem)};
  \node[box, right=of core] (pipe) {Pipeline\\Python / XGBoost};
  \node[box, right=of pipe] (rank) {Ranking\\JSON + SHAP};
  \node[box, right=of rank] (appdb) {BD operativa\\app};
  \node[box, below=1.1cm of appdb] (api) {API\\ASP.NET};
  \node[box, left=of api] (web) {React\\Prevencionistas};

  \draw[arr] (core) -- node[above, font=\tiny]{solo lectura} (pipe);
  \draw[arr] (pipe) -- (rank);
  \draw[arr] (rank) -- node[above, font=\tiny]{import} (appdb);
  \draw[arr] (appdb) -- (api);
  \draw[arr] (api) -- node[above, font=\tiny]{REST/JWT} (web);
\end{tikzpicture}

	\caption{Flujo de informaci\'on: core ART $\rightarrow$ pipeline $\rightarrow$ ranking $\rightarrow$ BD app $\rightarrow$ API $\rightarrow$ React.}
	\label{fig:flujo-info}
\end{figure}

\section{Arquitectura de componentes (desarrollo local)}

La Figura~\ref{fig:componentes-local} muestra el despliegue usado para la demo del 50\,\%.

\begin{figure}[htbp]
	\centering
	% Componentes / despliegue local
\begin{tikzpicture}[
  font=\scriptsize\sffamily,
  box/.style={draw=blue!70, fill=blue!6, rounded corners=2pt, align=center, minimum height=0.8cm},
  host/.style={draw=blue!40, dashed, inner sep=10pt},
  ]
  \node[box, minimum width=3.2cm] (browser) at (0,2.4) {Navegador\\React :5173};
  \node[box, minimum width=3.2cm] (api) at (0,1.0) {API .NET :5080};
  \node[box, minimum width=3.2cm] (sql) at (0,-0.4) {SQL Server\\PfiPrevencion};
  \node[box, minimum width=3.2cm] (ml) at (4.2,1.0) {Pipeline ML\\(offline)};
  \node[host, fit=(browser)(api)(sql), label={[blue!70]above:Estaci\'on local de desarrollo}] {};
  \draw[-{Latex}, thick, blue!60] (browser) -- node[right, font=\tiny]{/api proxy} (api);
  \draw[-{Latex}, thick, blue!60] (api) -- (sql);
  \draw[-{Latex}, thick, blue!60] (ml) -- node[above, font=\tiny]{import ranking} (sql);
\end{tikzpicture}

	\caption{Componentes en entorno local de demostraci\'on.}
	\label{fig:componentes-local}
\end{figure}

\section{Modelo de clases de dominio}

\begin{figure}[htbp]
	\centering
	% Diagrama de clases de dominio (simplificado)
\begin{tikzpicture}[
  font=\scriptsize\sffamily,
  cls/.style={
    rectangle split,
    rectangle split parts=2,
    rectangle split part fill={NavyBlue!14, white},
    draw=NavyBlue!70,
    rounded corners=3pt,
    align=left,
    text width=2.9cm,
    inner sep=4pt,
    line width=0.75pt,
    font=\scriptsize\sffamily,
  },
  rel/.style={line width=0.8pt, NavyBlue!68},
  mul/.style={font=\tiny\sffamily, NavyBlue},
]

  \node[font=\scriptsize\sffamily\bfseries, NavyBlue] at (0, 5.12) {Identidad};
  \node[font=\scriptsize\sffamily\bfseries, NavyBlue] at (4.6, 5.12) {Priorizaci\'on};
  \node[font=\scriptsize\sffamily\bfseries, NavyBlue] at (9.2, 5.12) {Operaci\'on};

  \node[cls] (u) at (0, 2.55) {\textbf{Usuario}
    \nodepart{second} Id\\Email\\Rol\\ZonaId};
  \node[cls] (s) at (4.6, 2.55) {\textbf{RankingSnapshot}
    \nodepart{second} Periodo\\Horizonte\\EsVigente\\ArtifactHash};
  \node[cls] (a) at (9.2, 2.55) {\textbf{AgendaItem}
    \nodepart{second} FechaPlanificada\\Estado\\MotivoVisita};

  \node[cls] (z) at (0, 0.15) {\textbf{Zona}
    \nodepart{second} Id\\Nombre};
  \node[cls] (i) at (4.6, 0.15) {\textbf{RankingItem}
    \nodepart{second} Ranking\\ScoreRiesgo\\Banda\\Motivo};
  \node[cls] (v) at (9.2, 0.15) {\textbf{VisitaApp}
    \nodepart{second} FechaRealizada\\Resultado};

  \node[cls] (c) at (4.6, -2.15) {\textbf{ContribucionSHAP}
    \nodepart{second} Feature\\ShapValue\\Direccion};

  \draw[rel] (u.south) -- (z.north)
    node[mul, left, pos=0.18] {*}
    node[mul, left, pos=0.82] {1};
  \draw[rel] (s.south) -- (i.north)
    node[mul, left, pos=0.18] {1}
    node[mul, left, pos=0.82] {*};
  \draw[rel] (i.south) -- (c.north)
    node[mul, left, pos=0.15] {1}
    node[mul, left, pos=0.85] {*};
  \draw[rel] (a.south) -- (v.north)
    node[mul, right, pos=0.18] {1}
    node[mul, right, pos=0.82] {0..1};
  % Snapshot $\to$ Agenda: misma cota Y (cajas de distinta altura)
  \draw[rel]
    ([yshift=-10pt]s.east) -- ([yshift=-10pt]s.east -| a.west)
    node[mul, above, pos=0.16] {1}
    node[mul, above, pos=0.84] {*};

  % Usuario $\to$ Agenda: tramo superior a Y fija, luego baja en vertical
  \coordinate (barra-ua) at ([yshift=0.88cm]u.north);
  \draw[rel, rounded corners=3.5pt]
    (u.north) -- (barra-ua) -- (barra-ua -| a.north) -- (a.north);
  \node[mul] at ($(u.north)+(-0.28,0.18)$) {1};
  \node[mul] at ($(a.north)+(0.28,0.18)$) {*};

\end{tikzpicture}

	\caption{Diagrama de clases simplificado del dominio de la aplicaci\'on.}
	\label{fig:clases}
\end{figure}

\section{Modelo de datos (ER)}

La BD operativa `PfiPrevencion` (schema \texttt{app}) almacena snapshots de ranking, \'items, contribuciones SHAP, outcomes, m\'etricas, coberturas, agenda y visitas. La Figura~\ref{fig:er} resume las entidades principales.

\begin{figure}[htbp]
	\centering
	% Modelo entidad-relaci\'on (crow's foot simplificado)
\begin{tikzpicture}[
  font=\tiny\sffamily,
  ent/.style={draw=blue!70, fill=blue!5, rounded corners=1pt, align=center, minimum width=2.1cm, minimum height=0.7cm},
  ]
  \node[ent] (zonas) at (0,3) {Zonas};
  \node[ent] (usuarios) at (3,3) {Usuarios};
  \node[ent] (snap) at (0,1.4) {RankingSnapshots};
  \node[ent] (items) at (3.5,1.4) {RankingItems};
  \node[ent] (shap) at (7,1.4) {Contribuciones};
  \node[ent] (out) at (3.5,-0.2) {Outcomes};
  \node[ent] (agenda) at (0,-0.2) {AgendaItems};
  \node[ent] (visitas) at (0,-1.8) {VisitasApp};
  \node[ent] (met) at (7,-0.2) {MetricasObs.};
  \node[ent] (cob) at (7,-1.8) {Coberturas};

  \draw[-{Latex}, blue!60] (usuarios) -- node[above, font=\tiny]{N:1} (zonas);
  \draw[-{Latex}, blue!60] (items) -- node[above, font=\tiny]{N:1} (snap);
  \draw[-{Latex}, blue!60] (shap) -- (items);
  \draw[-{Latex}, blue!60] (out) -- (items);
  \draw[-{Latex}, blue!60] (agenda) -- (usuarios);
  \draw[-{Latex}, blue!60] (agenda) -- (snap);
  \draw[-{Latex}, blue!60] (visitas) -- (agenda);
  \draw[-{Latex}, blue!60] (met) -- (snap);
  \draw[-{Latex}, blue!60] (cob) -- (snap);
\end{tikzpicture}

	\caption{Diagrama entidad-relaci\'on simplificado de la BD operativa.}
	\label{fig:er}
\end{figure}

Arquitectura de datasources:
\begin{enumerate}
	\item \textbf{BD core de la ART (on-prem):} siniestros, n\'omina, prevenci\'on, establecimientos. Solo la consume el pipeline.
	\item \textbf{BD operativa de la app:} rankings importados y operaci\'on diaria.
	\item \textbf{Artefactos ML:} modelos y reportes generados offline (no se cargan en cada request).
\end{enumerate}

\section{Arquitectura AWS Free Tier (objetivo)}

Para el hito 50\,\% la demo es local. La Figura~\ref{fig:aws-ft} documenta el despliegue objetivo posterior, alineado a AWS Free Tier personal (hasta ~USD\,200 / 6 meses): una EC2 con Nginx + React + API + SQL Express, S3 opcional para rankings y Budgets para control de cr\'editos. La BD core de la ART no se expone a Internet; el dise\~no contempla multi-ART.

\begin{figure}[htbp]
	\centering
	% Arquitectura AWS Free Tier objetivo (post-despliegue)
\begin{tikzpicture}[
  font=\scriptsize\sffamily,
  box/.style={draw=blue!70, fill=blue!6, rounded corners=2pt, align=center, minimum height=0.75cm},
  cloud/.style={draw=blue!45, dashed, inner sep=10pt},
  ]
  \node[box, minimum width=2.4cm] (user) at (0,3.2) {Usuario ART};
  \node[box, minimum width=2.6cm] (sg) at (0,2.0) {SG :80/:443};
  \node[box, minimum width=3.4cm] (nginx) at (-2.2,0.6) {Nginx\\React dist};
  \node[box, minimum width=3.4cm] (api) at (2.2,0.6) {API .NET};
  \node[box, minimum width=3.4cm] (sql) at (0,-0.8) {SQL Express\\(misma EC2)};
  \node[box, minimum width=2.6cm] (s3) at (5.2,0.6) {S3\\rankings};
  \node[box, minimum width=2.6cm] (bud) at (5.2,2.0) {Budgets};
  \node[box, minimum width=2.8cm] (core) at (-5.0,0.6) {BD core ART\\(on-prem)};
  \node[box, minimum width=2.8cm] (pipe) at (-5.0,-0.8) {Pipeline\\local/VPN};

  \node[cloud, fit=(sg)(nginx)(api)(sql), label={[blue!70]above:EC2 Free Tier (t3.micro/small)}] {};

  \draw[-{Latex}, thick, blue!60] (user) -- (sg);
  \draw[-{Latex}, thick, blue!60] (sg) -- (nginx);
  \draw[-{Latex}, thick, blue!60] (sg) -- (api);
  \draw[-{Latex}, thick, blue!60] (api) -- (sql);
  \draw[-{Latex}, thick, blue!60] (core) -- (pipe);
  \draw[-{Latex}, thick, blue!60] (pipe) -- (s3);
  \draw[-{Latex}, thick, blue!60] (s3) -- (api);
  \draw[dashed, blue!50] (bud) -- (sg);
\end{tikzpicture}

	\caption{Arquitectura objetivo en AWS Free Tier (EC2 todo-en-uno).}
	\label{fig:aws-ft}
\end{figure}

\section{Avance de implementaci\'on (demo)}

Al momento del hito 50\,\% se encuentra operativo un prototipo local con las pantallas de login, resumen, priorizaci\'on, detalle (explicaci\'on + SHAP), agenda, estad\'isticas, exportaciones y perfil. La API expone autenticaci\'on JWT y endpoints de riesgo, agenda y anal\'itica hist\'orica.

\subsection{Resultados del modelo (referencia holdout 2025)}

\begin{table}[htbp]
	\centering
	\caption{M\'etricas de referencia del modelo en holdout 2025.}
	\label{tab:metricas-holdout}
	\begin{tabular}{@{}lcc@{}}
		\hline
		\textbf{M\'etrica} & \textbf{1m} & \textbf{3m} \\
		\hline
		PR-AUC & 0,366 & 0,515 \\
		Precision @ top-50 / mes & 76,4\,\% & 93,5\,\% \\
		Precision @ top-225 / mes & 51,5\,\% & 79,0\,\% \\
		Visitas reales @ mismo volumen & 10,3\,\% & 17,8\,\% \\
		\hline
	\end{tabular}
\end{table}

Las pantallas del prototipo (Figuras~\ref{fig:mock-login}--\ref{fig:mock-estadisticas} del Cap\'itulo~\ref{ch:descripcion}) constituyen la evidencia de avance funcional exigida por la r\'ubrica. El recorrido de demostraci\'on en vivo es: login $\rightarrow$ resumen $\rightarrow$ riesgo $\rightarrow$ detalle/SHAP $\rightarrow$ agenda $\rightarrow$ estad\'isticas.

\section{Limitaciones del hito}

\begin{itemize}
	\item User research en curso: encuesta $n=81$ (Anexo~\ref{anexo:encuesta}) y primera entrevista (Anexo~\ref{anexo:entrevista-adrian-elias}); restan dos entrevistas y las \emph{user personas}.
	\item Despliegue cloud diferido a post-hito (misma arquitectura Free Tier documentada).
	\item El score de riesgo no debe interpretarse como probabilidad calibrada de accidente.
\end{itemize}

\end{contenidoNuevo}
